\section{2020--2025: cómo el ánimo de Silicon Valley pasó de los despidos al predominio de la IA}

Esta sección revisa primero cómo cambió el ánimo del mercado entre 2020 y 2025, porque los juicios de inversión de Freda no surgieron de la nada: se formaron a lo largo de varios ciclos sucesivos de tasas de interés, inflación, despidos, financiamiento de IA y volatilidad en el mercado secundario.

Este capítulo comienza por establecer el contexto temporal. Freda repasa las palabras clave de Silicon Valley entre 2020 y 2025: en 2020, la pandemia y los recortes de tasas; en 2021, un mercado que siguió en auge, aunque las empresas de pequeña y mediana capitalización tocaron techo primero; en 2022, la inflación y la fuerte caída del Nasdaq; en 2023, la IA se convirtió de pronto en el eje del mercado antes de que este alcanzara a reaccionar; 2024 fue el año en que la narrativa de la IA resultó más fácil de sostener, y en 2025 se sumaron los aranceles, la relación entre China y Estados Unidos y una ola de IA que siguió acumulando fuerza.

La observación más importante aquí es que las señales de financiamiento de la IA se adelantaron a la narrativa general. En 2022 el mercado secundario se veía sombrío, pero empresas de modelos como Inflection AI, Anthropic, Cohere y Stability todavía lograban levantar montos inusualmente grandes en sus rondas Serie A. La explicación de Freda es directa: Scaling Law, la compra de GPU y el gasto de capital fueron conceptos que hicieron que el dinero de los VC fluyera muy pronto hacia Nvidia y hacia las empresas de modelos. Aquí aparece la ventaja de los fondos Crossover: si por el momento no se puede saber cuál de las decenas de empresas de modelos saldrá ganadora, se puede expresar una postura sobre la demanda total de IA a través de Nvidia en el mercado secundario.

\begin{knowledgebox}{Términos clave: Crossover, Beta y Alpha}
\begin{tabularx}{\textwidth}{p{0.22\textwidth}p{0.32\textwidth}X}
\toprule
Término & Explicación breve & Papel en este episodio \\
\midrule
Fondo Crossover & Fondo que invierte a la vez en empresas privadas del mercado primario y en empresas cotizadas del mercado secundario & Puede usar el mercado público (Nvidia, Meta, Robinhood, entre otras) para expresar su juicio sobre las tendencias de frontera. \\
Beta & Rendimiento que proviene del alza general del mercado o del sector & Si solo se apuesta por el ciclo de las criptomonedas, comprar Coinbase no necesariamente aporta más Alpha que comprar bitcoin. \\
Alpha & Rendimiento propio de la empresa que va más allá de lo que explican el sector o el índice & Robinhood sirve para mostrar cómo la capacidad de ejecución de una empresa y la revaloración de su negocio generan Alpha. \\
CAPEX & Capital Expenditure, gasto de capital & Los centros de datos de IA, las GPU, el entrenamiento y la infraestructura en la nube dependen en gran medida del CAPEX. \\
\bottomrule
\end{tabularx}
\end{knowledgebox}

\subsection{Resumen de la sección}

Entre 2020 y 2025, Silicon Valley no tuvo una recuperación gradual, sino una serie de cambios bruscos y continuos. Antes de que la IA se volviera el eje principal, el capital ya había detectado señales inusuales en las grandes rondas de financiamiento de las empresas de modelos y en la demanda de Nvidia. Solo con este contexto se entiende por qué este episodio pasa constantemente del mercado primario al secundario y viceversa.
